\chapter{A Compatible Smooth Structure and the Topological Endpoint}
\label{ch:smoothing}

The Ricci-flow argument is formulated for smooth manifolds.  The original
Poincare question is topological.  The bridge between these statements is a
construction: a closed topological three-manifold is given a smooth model,
with one specified homeomorphism relating the two carriers.  This chapter
follows the construction used by the formalization.  It is useful to keep
the logical order visible.  Hamilton's relative handle argument constructs
the compatible PL charts; Cairns's construction smooths a finite
triangulation.  Neither step assumes simple connectedness.  Only after the
model has been made smooth do we transport simple connectedness and invoke
the closed-manifold topology package used by the smooth endpoint.

Throughout, $M$ is a topological manifold with a Hausdorff
topology, a second-countable topology, compactness, and a (possibly merely
topological) charted space over $E_3=\mathbb R^3$.  Connectedness and
nonemptiness are hypotheses of the smoothing input.  A smooth model is a
manifold $N$ with its own topology and $C^\infty$ atlas;
it is not an assertion that the initially supplied charts on $M$
are smooth.  We write $h:M\cong N$ for the chosen homeomorphism.

\section{A compatible smoothing}
\label{sec:smoothing-contract}

\begin{definition}[Compatible smooth model]
\label{def:smoothing-bridge}
For a connected nonempty topological manifold $M$, a compatible smooth
model consists of
\begin{enumerate}
\item a manifold $N$ with a topology and an atlas over $E_3$;
\item a $C^\infty$ manifold structure on $N$;
\item a Hausdorff, second-countable and compact topology on $N$;
\item connectedness of $N$; and
\item a specified homeomorphism $h:M\cong N$.
\end{enumerate}
The smooth atlas must induce this same topology on $N$.
\end{definition}

The homeomorphism is part of the data because all later transport must use
the same map.  In particular, replacing $h$ by an existentially chosen
homeomorphism after applying the smooth theorem would lose the relation
between the transported topological hypotheses and the endpoint map.

\begin{theorem}[Hamilton--Cairns bridge]
\label{thm:m76}
For every compact, Hausdorff, second-countable, charted topological
three-manifold that is nonempty and connected, there is a compatible
smooth model. The proof is independent of simple connectedness.
\lean{PoincareMT.m76CompatibleSmoothing}
\lean{PoincareMT.M76.smoothingConclusion_of_lower_handle_cases}
\source{Hamilton1976,Cairns1940}
\dcref{Hamilton 1976, Theorem 2(1), pp.~64,69; Cairns 1940, Theorem III, p.~797 and Section 11, p.~807}
\end{theorem}

The rest of the chapter proves the content hidden behind this theorem.  The
four chart-handle indices are represented by finite subsets
$J\subseteq\{0,1,2\}$.  A straightening for $J$ is required for every
coordinate chart $e:E_3\dashrightarrow E_3$ whose source contains the
coordinate cylinder
\[
 C_J=\{x\in E_3: |x_i|\leq 1\text{ for }i\in J\}.
\]
For every open neighborhood $N$ of the frontier of $C_J$, and every
map $e$ that is locally piecewise affine on $\operatorname{dom}(e)\cap N$, the conclusion
is a homeomorphism $A:E_3\cong E_3$ for which $e\circ A$ is finite
piecewise affine on the closed unit ball, together with a homotopy from the
identity to $A$.
Every homotopy stage is a homeomorphism, is the identity outside the
radius-two ball, and fixes $C_J^c\cup\partial C_J$ pointwise.
The correction is supported in the prescribed cylinder.

\section{The protected end and Wall cores}
\label{sec:wall-core}

The lower-index argument works in a noncompact PL domain cut out of a chart.
The end condition is stronger and more relative than saying that one
complement is simply connected.

\begin{definition}[One simply connected end]
\label{def:one-end}
For a topological space $Y$, write $Y$ has one simply connected end when
for every compact $C\subset Y$ there is a compact $D\subset Y$ such that
\[
C\subset\operatorname{int}D,\qquad D^c\text{ is connected},
\]
with $D^c$ nonempty, and every path $p:[0,1]\to Y$ with $p(0)=p(1)=x$
and image disjoint from $D$ admits a based nullhomotopy avoiding $C$.
The homotopy is not
required to avoid $D$; the protected set is the smaller $C$.
\lean{PoincareMT.M76.HasOneSimplyConnectedEnd}
\source{Hamilton1976}
\dcref{Hamilton Lemma 2, p.~65}
\end{definition}

Let $e=(e_i)$ be an ambient PL atlas on a Hausdorff paracompact space $X$.
A PL domain has the compatible half-space charts of
Definition~\ref{def:compact-pl-domain}; connectedness and compact frontier
are additional hypotheses. Wall's relative conclusion is the
following exact statement.

\begin{proposition}[Protected Wall compact core]
\label{prop:wall-core}
Assume $R$ is a connected PL domain, $\partial R$ is compact, $R$ has one
simply connected end, and $A\subset R$ is compact.  There are compact
$K\subset R$ and a set $S\subset\operatorname{int}R$ such that
\begin{align*}
 &K\text{ is a PL domain},\quad S\text{ is a nonempty chartwise PL sphere},\\
 &S\cap\partial R=\varnothing,\qquad
   \partial K=\partial R\cup S,\\
 &A\cup\partial R\subset\operatorname{int}_R K,\qquad
   \partial_R K=S.
\end{align*}
The last equality is a relative-frontier equality in the subtype $R$.
No connectedness of $K$ is part of the exported predicate, although the
construction below produces a connected core.
\lean{PoincareMT.M76.hasWallCompactCore}
\lean{PoincareMT.M76.HasWallCompactCore}
\source{Hamilton1976}
\dcref{Hamilton Lemma 2, pp.~65--66}
\end{proposition}

\begin{proof}
Apply Theorem~\ref{thm:protected-wall-core}. Its minimal-complexity
compression and marked-disk contradiction first give a finite spherical
frontier. Filling the bounded complementary components leaves one connected
exterior; a two-foot ball along an exterior arc joins two frontier spheres.
Refilling the bounded components preserves the protected relative interior
and decreases the sphere count. These are precisely the ambient and
relative frontier identities retained by that theorem.
\end{proof}

\section{Marked Dehn surfaces and protected prime replacement}
\label{sec:dehn-prime}

Wall cores alone do not straighten a lattice handle.  The lower-index
models have transverse lattice $L$ and a protected compact region.  A
generalized Dehn input supplies, for every relevant chart, marked proper
annuli and disks whose boundaries are fixed on the protected collar.  The
boundary parametrizations are part of the objects: an unmarked embedded
surface is insufficient because the next projection must agree on the old
frontier.

\begin{definition}[Generalized Dehn input]
\label{def:generalized-dehn}
The generalized Dehn input is the conjunction of
\begin{enumerate}
\item the protected index-one annulus, with both prescribed rim maps;
\item the protected index-two pair of simultaneously disjoint disks,
with their prescribed negative and positive rim maps; and
\item the standard proper-disk theorem for every compact standard PL
domain, prescribed finite PL boundary circle and continuous filling.
\end{enumerate}
The first two clauses quantify over the supplied full discrete lattice
and atlas, with the original chart and retained-block hypotheses of
Section~\ref{sec:dehn-protected-disks}.  Their surfaces avoid the
protected block, are proper against the whole old frontier, and come
with the enclosing region formed from those same surfaces and the full
attaching patch.  The third clause preserves every value of the
prescribed circle and keeps the disk interior off the whole domain
frontier, as in Theorem~\ref{thm:dehn-standard-proper}.
\lean{PoincareMT.M76.HasHamiltonGeneralizedDehnInput}
\lean{PoincareMT.M76.hasHamiltonGeneralizedDehnInput}
\source{Hamilton1976}
\dcref{Hamilton 1976, pp.~65--69}
\end{definition}

The finite constructions are proved in
Sections~\ref{sec:dehn-marked-tower} and
\ref{sec:dehn-annulus-descent}.  Tower ascent decreases
$|T|-|s_T(T)|$ for one fixed finite set of source marks $T$.
During disk descent the number of boundary double components strictly
decreases without increasing the interior count; after the boundary
count vanishes, interior circle surgery decreases the latter count
while fixing the entire rim.  Annulus descent instead decreases the
number of connected components of the actual source double locus,
with both whole rim maps fixed.  Finite supported repairs and exact
tube preimages justify these component comparisons.  The separate
proper-disk boundary correction uses the finite residual
\eqref{eq:dehn-collar-residual} to paste its collar compression to
the identity on the rest of the domain.  That residual construction
is part of the proof, rather than an additional surface-existence
premise in the generalized Dehn input.

For a lattice $L$ and a chart $e$, the protected irreducible replacement
predicate is an atlas-level relative conclusion.  Given a marked protected
ball $D$, it produces a chart family on the same lattice-handle ambient
space and an open neighborhood of $D$ together with the old frontier.  The
new family is PL-irreducible on the handle, and the identity map is
chartwise PL in both directions between the old and new families on the
corresponding preimage.  Thus the predicate records the actual relative
replacement data; it does not assert a homotopy equivalence to a product
model. Theorem~\ref{thm:protected-prime-replacement} constructs it by
maximal sphere cuts, relative capping and retained-chart transport. Index
$i=0,1,2$ counts interval coordinates; the transverse torus dimensions are
respectively $3,2,1$. Thus the cases are a closed three-torus, an
interval times a two-torus, and a solid torus. Relative torus rigidity is
applied after this replacement, to the resulting irreducible structure.

\section{The four handle cases}
\label{sec:handle-cases}

For indices zero, one and two, the preceding relative-rigidity chapter
combines these inputs and constructs the required supported correction.
In particular, the comparison $A$ precedes the placement $Q$; the final
correction is $B=A^{-1}Q$. The inclusion $Q(C)\subset A(P)$ puts its
inverse in the protected agreement neighborhood, where both PL identity
directions are available. The proof there establishes finite PL regularity
of the corrected original chart, as well as an isotopy fixing the entire
cylinder frontier and the radius-two exterior.

The fourth case, $J=\{0,1,2\}$, uses Alexander's finite height induction
(1924, pp.~6--8). Here the coordinate cylinder is the closed unit ball
$B=[-1,1]^3$ for the maximum norm. Write $h$ for the given chart and
$S=h(\partial B)$. The chart is PL near $\partial B$; compactness therefore
gives a finite PL parameterization of this sphere. Choose a compact convex
polyhedron $C$ with $h(B)\subset\operatorname{int}C$, and put
$Z=\partial(C\times[-1,1])$. We prove the stronger assertion that there
is an open connected $U$, with $\partial U=S$ and
$\overline U\subset\operatorname{int}C$, for which both
\[
 (\overline U,S),\qquad
 \bigl(Z\setminus(U\times\{1\}),S\times\{1\}\bigr)
\]
are finite PL three-ball pairs. The second ball is the complementary side
in the fixed polyhedral three-sphere $Z$.

Triangulate $S$ and choose an affine height $a$ distinct on its finitely
many vertices. Each whole section $S\cap\{a=c\}$ has a finite presentation
by simple polygons and at most one residual point; distinct polygons can
meet only at that point. Give the section charge $q(c)=0$ if its polygons
are pairwise disjoint, and otherwise give it charge equal to the number of
\emph{all} polygons in the presentation. Only finitely many levels have
positive charge, so
\[
 N=\sum_{c\in\operatorname{supp}q}q(c)
\]
is a nonnegative integer. Recursive surfaces retain such complete section
presentations, finite height-event sets, collars of branching sections,
and approaches from both height directions at nonisolated section points.
We induct on $N$, keeping the same outer body $C$ and both ball pairs.

If $q(c)>0$, the planar polygon arrangement at height $c$ supplies a
polygon bounding a planar disk $d$ whose intersection with $S$ is exactly
its rim. Cutting the parameter sphere at that rim gives two surface disks
$s,s'$ with $S=s\cup s'$ and $s\cap s'=\partial d$; each capped sphere
is $s\cup d$ or $s'\cup d$. Choose a small height window avoiding the other
finitely many events. The branching collars push the two caps to opposite
sides inside this window by finite PL ambient homeomorphisms supported in
$\operatorname{int}C$. When the common branching point lies on the chosen
rim it is retained by the two pointed cap moves; otherwise ordinary cap
moves introduce only isolated extreme sections. Whole section polygons
are transported or partitioned between the children. Their charges obey
$q_1(t)+q_2(t)\leq q(t)$ at every corresponding level, with a strict
inequality at the selected level; new isolated extreme sections have
charge zero. Consequently $N_1+N_2<N$. This is the finite decrease needed
for both recursive calls, without requiring a decrease in vertex count.

Pull the resulting regions and both of their ball pairs back by the same
cap homeomorphisms. The actual bounded regions $U_1,U_2$ of the capped
spheres have three possible incidences: their closures meet exactly in
$d$, or one open region lies inside the other. In the first case the new
region is $\operatorname{int}(\overline U_1\cup\overline U_2)$; attachment
of two balls along their common boundary disk gives its ball. In the
nested case, say $U_1\subset U_2$, the region is
$U_2\setminus\overline U_1$. To obtain its ball, take a collar of the
appropriate surface disk inside the certified exterior of $U_1$, avoiding
the exterior of $U_2$ except at the common rim. The spherical covering
forces this collar into $\overline U_2$. Ball attachment followed by
compression across the collar removes $\overline U_1$ and leaves a ball
with boundary $s\cup s'$. The exterior certificates provide the relative
neighborhood along the old cap needed to glue this compression to the
identity. The complementary ball is obtained by the corresponding
excision in the first case, and in the nested case by attaching
$\overline U_1\times\{1\}$ to the exterior of $U_2$. These operations
identify the literal set $Z\setminus(U\times\{1\})$, as well as its
boundary, and so preserve the strengthened induction assertion.

For the base $N=0$, zero charge initially says only that each section's
polygons are disjoint. Use the two-sided approach property to tilt the
height slightly on the same finite surface complex so that its positive
and negative vertex links are connected whenever nonempty. The generic
height has unique extreme vertices $p_-,p_+$; connectedness of the sphere
and the signed links then imply that every strict intermediate section
is one whole simple polygon. Let $d(c)$ be its planar disk. Near $p_-$ the
surface is conical, giving a first three-ball bounded by $d(c)$ and the
whole lower surface. Its exterior is also a ball: in a small convex
envelope the complementary frontier disk cones to a second ball, and
attaching that cone to the exterior of the envelope in $Z$ gives precisely
the exterior of the first cone. Reversing height gives the corresponding
last cone near $p_+$.

Between these cones, a collar of the whole section gives one joint finite
PL cylinder chart for a small height band. In its height coordinates a
small move from $c$ to $c'$ is realized by
$(x,t)\mapsto(x,t+(c'-c)\sigma(x,t))$, where the PL cutoff $\sigma$ is one
on the swept disk and vanishes outside a compact subset of the chart.
Taking $|c'-c|$ small makes each vertical map strictly increasing. After
conjugation this is an ambient homeomorphism supported inside $C$ that
carries both the capped lower ball and its exterior to the next ones.
The cylinder slices equal the previously chosen $d(c)$ by uniqueness of
the planar disk with its specified polygon rim. A finite subdivision of
the compact interval between the two endpoint heights, subordinate to
these collar neighborhoods, propagates both ball pairs to the last cone.
Attaching that cone removes the final cap and gives the two required
balls for the original sphere. Thus the base and the positive-charge
reduction together prove the strengthened assertion.

It remains to identify the constructed region with the given chart image.
Set $V=h(\operatorname{int}B)$. The chart homeomorphism and compactness give
$\partial V=S$ and $\overline V=h(B)$. If $U$ and $V$ were disjoint, the
connected interior of the certified exterior ball,
$Z\setminus(\overline U\times\{1\})$, would meet $V\times\{1\}$ and
avoid its relative frontier $S\times\{1\}$, hence would lie entirely in
$V\times\{1\}$. But it contains a point of the bottom face
$C\times\{-1\}$, a contradiction. Since $U$ and $V$ intersect, their
connectedness and common frontier imply both inclusions, so $U=V$.
The finite PL ball is therefore the actual pair $(h(B),h(\partial B))$.

Extend the prescribed finite PL boundary map to a finite PL homeomorphism
$g:B\to h(B)$ agreeing with $h$ on $\partial B$, by coning boundary maps
in ball parameterizations. Then $A=h^{-1}g$ fixes $\partial B$ and extends
by the identity outside $B$, with $hA=g$ on $B$. The corrected chart is
PL near the whole ball: on the ball it is $g$, outside it retains $h$,
and these PL maps agree along the polyhedral boundary. For $t>0$ the
Alexander isotopy is $A_t(x)=tA(x/t)$ for $\|x\|\leq t$, and $A_t(x)=x$
otherwise. It starts at $A_0=\operatorname{id}$ and fixes the entire unit
sphere and exterior at
every time; the same formula for the inverse gives joint continuity of
the inverse family. This completes index three. Together with indices
zero, one and two it gives the straightening for every
$J\subseteq\{0,1,2\}$ with the hypotheses and quantifiers of
Section~\ref{sec:smoothing-contract}.

\section{Finite supported assembly}
\label{sec:finite-assembly}

The local straightening theorem is converted into a global supported
correction as follows.  Choose a finite compact chart presentation of a
compact set $Q$ in the old PL atlas.  Order its finitely many affine faces
by a Hamilton face order.  If $P_k$ is the union of the first $k$ model
faces, then
\[
 P_0=\varnothing,\qquad P_{k+1}=P_k\cup\operatorname{carrier}(\sigma_k),
 \qquad Q\subset P_n.
\]
At stage $k$, the new face frontier lies in the current open PL
neighborhood $U_k$ of $P_k$, and its intrinsic interior is disjoint from
$P_k$. Choose the handle width after $U_k$ has been chosen: its attaching
region then lies in $U_k$, while its moving interior avoids $P_k$.
Straightening this handle extends the PL region to a neighborhood of
the new face and $P_k$. The correction fixes $P_k$, is supported on a
compact subset of the chart domain, and extends by the identity outside
that domain. Compose it with the accumulated correction. The new open
PL neighborhood can shrink around the old carrier, but must contain
the enlarged carrier $P_{k+1}$. Thus the induction preserves the
previously treated faces without presupposing a single handle width
that works for every later neighborhood.

The result is an ambient homeomorphism $F$, an open set $U$ containing $Q$,
and a compact support set $S$ such that $F=\mathrm{id}$ on $S^c$. If
$d$ is the incoming chart and $c_i$ an old chart, then
$(d\circ F)\circ c_i^{-1}$ is locally piecewise affine on
$c_i(U\cap\operatorname{dom}c_i)$. This construction applies on the
actual overlap, where $d$ has full source.

To obtain an atlas, choose finitely many compact chart cores whose
interiors cover $M$. Suppose a finite compatible family already covers
their preceding union $A$, and the next core is $B\subset\operatorname{dom}d$.
Apply the construction to the compact overlap $A\cap B$. Its compact
support lies inside the actual overlap, so its homeomorphism extends by
the identity to $M$. The corrected chart still has the same source and
target as $d$. Choose open neighborhoods $U_A,U_B$ of $A,B$ with
$U_A\cap U_B$ in the straightened region. Restrict the old charts to
$U_A$ and the corrected chart to $U_B$. Old transitions remain PL, and
the displayed formula proves this for the new transitions; their inverses
are PL because they are local homeomorphisms. The new family covers
$A\cup B$. Finite induction constructs a compatible PL atlas on the
original topology, with the cocycle identities inherited from its actual
coordinate maps.

\section{Independent vertices before Cairns smoothing}
\label{sec:independent-triangulation}

Here is the finite geometric realization of the compatible PL atlas just
constructed. Choose finitely many charts $e_i$ and compact cores
$Q_i\subset\operatorname{dom}e_i$ whose interiors cover $M$. In each
coordinate domain, enclose $e_i(Q_i)$ in an inner finite polyhedron and
then in the interior of an outer finite polyhedron, still inside that
domain. On a common subdivision assign the value $(1,v)$ to each vertex
$v$ in the inner polyhedron and zero to the remaining vertices, and
extend affinely over simplices. This gives a PL graph block equal to
$x\mapsto(1,x)$ on the inner polyhedron and zero on the outer boundary;
extend it by zero outside. Its height lies in $[0,1]$. If the height is
one, every vertex with positive barycentric weight has height one, so
the second coordinate is exactly $x$. Transporting the block by $e_i$
and extending by zero gives a continuous map
$g_i:M\to\mathbb R\times\mathbb R^3$, supported compactly inside the
chart. It is locally PL in every chart by compatibility of transitions.
The map
\[
 \Gamma:M\longrightarrow E=\prod_i(\mathbb R\times\mathbb R^3),
 \qquad \Gamma(x)=(g_i(x))_i
\]
is injective: choose $i$ with $x\in Q_i$; equality of the $i$th blocks
forces $y$ into the same chart with $e_i(y)=e_i(x)$. Compactness and the
Hausdorff property make $\Gamma$ a homeomorphism onto its image.

For each point, a finite local PL witness in a coordinate domain gives
a finite geometric complex in $E$: its affine image is embedded because
$\Gamma$ is injective. Choose the witnesses with the point in the
interior of their coordinate carriers. These interiors cover $M$, so
finitely many image complexes cover exactly $\Gamma(M)$. To triangulate
their union, describe their simplices by finite affine inequalities,
subdivide a surrounding polyhedral neighborhood along all the associated
hyperplanes, and retain the faces lying in at least one supplied simplex.
This is a finite geometric complex $K$ with $|K|=\Gamma(M)$; intersections
have been subdivided compatibly, and no new points are added to the image.

On $\Gamma(\operatorname{int}Q_i)$ the ambient linear projection $a_i$
onto the second coordinate of block $i$ equals $e_i\circ\Gamma^{-1}$.
Repeated barycentric subdivision makes every whole closed vertex star
lie in one of these open cores: the mesh tends to zero, the diameter of
a star is at most twice the mesh, and the finite core cover has a
Lebesgue number. Keep the name $K$ for the resulting subdivision.
Thus $a_i$ is injective on that star. A closed vertex star is a
neighborhood of its vertex in $|K|$, and the corresponding chart is
open, so $a_i$ sends the vertex into the interior of the star image in
$\mathbb R^3$.

These affine star maps also give the exact incidence conditions needed
below. For a face $\sigma$, choose one of its vertices and the associated
star map. The image of $\sigma$ has relative-interior points in the
ambient interior of this star image. Nearby points off the finitely many
lower-dimensional faces lie in tetrahedra; taking a limit and using the
complex intersection law gives a tetrahedral coface of $\sigma$. For an
edge, take a small ball about such a relative-interior point inside the
closed edge star; finiteness lets us avoid every simplex not incident
there. The ball minus the edge's affine line is connected. Normalizing
the barycentric weights outside the edge defines a continuous projection
onto its entire link: every
link point is reached along a sufficiently short segment from the chosen
point. The edge link is therefore connected. For a triangle, the two
sides of a small ball give two tetrahedral cofaces; a third would overlap
one of them in its interior, contrary to the complex intersection law.
Its link consequently has exactly two vertices. These are statements
about the full links of $K$, since the closed star of a vertex of
$\sigma$ contains all cofaces of $\sigma$. The construction has supplied
the affine star and link hypotheses, but its geometric vertices need
not be globally affinely independent.
\source{Hamilton1976,Cairns1940}
\dcref{Hamilton p.~69; Cairns pp.~798--799 and 804--807}

This issue is repaired before applying Cairns.  Since $K$ has finitely many
vertices, choose the Euclidean space $F=\mathbb R^{V(K)}$ and the basis
vector assignment $v\mapsto e_v$.  The affine extension theorem constructs
a map $f:E\to F$ satisfying
\[
 f|_\sigma\text{ affine for every face }\sigma,
 \qquad f|_{|K|}\text{ injective},
 \qquad f(v)=e_v.
\]
The basis vectors are affinely independent on every face.  A linear
synthesis map $Q:F\to E$ is a left inverse on $|K|$; hence injectivity of
$f$ is proved, rather than assumed.  The image complex $f(K)$ has exactly
the same face and link data as $K$, and the homeomorphism to the carrier is
composed with this embedded image.
\lean{Geometry.SimplicialComplex.exists_independent_euclidean_realization}
\lean{PoincareMT.M76.smoothingConclusion_of_finite_geometric_triangulation}
\source{Cairns1940}
\dcref{Cairns Section 3, pp.~798--799}

This re-realization is logically prior to the smoothing theorem.  It is not
an optional generic-position perturbation: it preserves the labelled faces,
their links, the carrier homeomorphism, and the affine star hypotheses while
providing the independent vertex set required by the Cairns construction.

\section{Cairns's atlas and pullback}
\label{sec:cairns-pullback}

For a finite independent three-complex $L\subset F$ with pure four-vertex
cofaces, connected links at edges, and two vertices in each triangle link,
construct an ambient plane field $P(x)$ of codimension three as follows.
At a vertex, affine independence extends the simplexwise affine
Brouwer-star map to an ambient affine map. Its linear part is surjective
because the star image has interior in $\mathbb R^3$, and its kernel is
uniformly transverse to the star's secants. This supplies an initial
plane. Here secant transversality means that projection with this kernel
has a positive lower bound on the norm of every unit secant of the
closed face star, not just on the directions of its individual edges.

The extension over positive-dimensional faces uses the special link
dimensions in dimension three. An edge link is a connected finite graph
with degree two at every vertex, hence a cycle. After quotienting by the
edge direction, fix two adjacent normal rays at $1,i\in\mathbb C$.
Radial normalization contracts all other positive radii to one; the
remaining angular gaps lie in $(0,\pi)$, have sum $2\pi$, and have
first gap $\pi/2$. These conditions form a nonempty convex set, so the
normalized cycle-projection space is contractible. A triangle has two
normal rays $u,v$; its normalized operators satisfy $Qu=1$ and $Qv<0$,
again a nonempty convex set. A tetrahedron has empty link and only the
zero-dimensional normal model. Restoring the face tangent directions
and the unrestricted operators on the ambient complement preserves
contractibility. Thus for each positive-dimensional face the entire
space of codimension-three planes transverse to its closed star is
contractible \cite[Sections 4--5 and 11, pp.~800--802, 807]{Cairns1940}.

Process the finitely many faces in increasing dimension, keeping the
field unchanged on a neighborhood of the processed carrier. A full
tetrahedral coface supplies an affine slice $b:\mathbb R^3\to F$ with
linear part $J$; transverse planes are represented by normalized
operators $Q:F\to\mathbb R^3$ with $QJ=1$ and kernel equal to the
plane. On a new face, choose an inner convex core whose surrounding
collar lies in the old field domain. Contractibility fills its boundary
operator values across the core. Retain the additional slice coordinates
as parameters, smooth relative to a neighborhood of the boundary, and
normalize back to $QJ=1$. Transversality survives because its operator
target is open and the controlled core is compact. This gives a smooth
full-slice field $q(u)$ agreeing with the old one near that boundary.
To extend it into the ambient space, fix $u_0$ and use
\[
 \Phi(u,v)=b(u)+v-Jq(u)v,\qquad v\in\ker q(u_0).
\]
On a compact convex base neighborhood, $q$ has a Lipschitz constant
$M$. Equality of two values of $\Phi$ first forces equality of their
fixed-kernel coordinates $v$, then gives
$u-u'=(q(u)-q(u'))v$. Thus $M\|v\|<1$ excludes a collision unless
$u=u'$. Along the zero section the derivative is invertible;
restricting this uniformly injective tube to the open invertibility
locus gives a smooth inverse neighborhood of the compact core. Set the new
ambient operator at $\Phi(u,v)$ equal to $q(u)$. Its kernel is constant
locally along the affine leaves parametrized by $v$. Equality of the
full-slice germs forces equality of these ambient leaf fields near the
old boundary, so they paste while preserving the previous field on a
neighborhood. The finite induction yields $P$ on an open neighborhood
of $|L|$, with $\dim P(x)+3=\dim F$, transverse to each required
face star and locally constant along its affine leaves
\cite[Sections 7--8, pp.~803--805]{Cairns1940}.

For the resulting charts, take $a$ in the relative interior of a face
and choose a full-coface frame $J:\mathbb R^3\to F$. On a smaller
ambient neighborhood there is a smooth normalized representative
$Q(y)$ with $Q(y)J=1$ and $\ker Q(y)=P(y)$. Define
$c(y)=Q(y)(y-a)$ on the carrier. The fixed projection $Q(a)$ is
injective on the closed face star with a bounded inverse on secants.
The paired triangular cofaces ensure that its image contains a
neighborhood of $Q(a)a$. The ambient formula $y\mapsto Q(y)(y-a)$
has derivative $Q(a)$ at $a$: the derivative of $Q$ is multiplied
by $a-a=0$. After composing with the inverse of the centered fixed
projection and shrinking the neighborhood, $c$ is a
perturbation of the identity whose error has Lipschitz constant less
than one. It is therefore a local homeomorphism onto an open subset
of $\mathbb R^3$. With $b(z)=a+Jz$, normalization and leaf constancy
give, on this chart,
\[
 b(c(y))-y\in P(y),\qquad Q(b(c(y)))=Q(y).
\]
Both points and their connecting segment lie in a sufficiently small
convex ambient field neighborhood, so the local leaf identity applies
along the whole segment.

Choose such charts $c_i$ covering $|L|$, with slices $b_i$ and linear
parts $J_i$. At $y=c_i^{-1}(z)$ in an overlap, both chart operators
have kernel $P(y)$ and satisfy $Q_j(y)J_j=1$. Consequently
$Q_i(b_i(z))J_j$ is invertible. The two slice points differ by a
vector in $P(y)$, so the overlap map is exactly
\[
 (c_j\circ c_i^{-1})(z)
 =\bigl(Q_i(b_i(z))J_j\bigr)^{-1}
   Q_i(b_i(z))\bigl(b_i(z)-b_j(0)\bigr).
\]
Every term on the right is smooth on the overlap, including inversion
on the open set of invertible matrices. This constructs the compatible
smooth atlas on the given carrier topology
\cite[Section 9, p.~806]{Cairns1940}.

If $e:M\cong|L|$ is the carrier homeomorphism and $a$ is the smooth
charted space on $|L|$, the smooth atlas on the original topological carrier
is
\[
e^*a.
\]
Explicitly, if $c_i$ are the smooth charts on $|L|$, the charts on $M$
are $\widetilde c_i=c_i\circ e$. Their transition identity is
\[
 \widetilde c_i\circ\widetilde c_j^{-1}
 =c_i\circ e\circ e^{-1}\circ c_j^{-1}=c_i\circ c_j^{-1}
\]
on the actual overlap, so it is smooth. Conversely, a smooth atlas on $M$
gives a compatible smooth model by taking $N=M$ and $h=\mathrm{id}$.
Thus
\[
 M\text{ has a compatible smooth model}\quad\Longleftrightarrow\quad
 \exists a\;\text{a smooth atlas on the original topology of }M.
\]
The forward implication uses the specified $h$ to pull back the model atlas;
it does not identify the original arbitrary charted space with that atlas.

\section{Transporting the topological hypotheses}
\label{sec:m77-m80}

Now assume in addition that $M$ is simply connected and choose a
compatible smooth model $(N,h)$. The inverse homeomorphism
$h^{-1}:N\cong M$ induces a homotopy equivalence,
so
\[
 \operatorname{SimplyConnected}(M)
 \Longrightarrow \operatorname{SimplyConnected}(N).
\]
The same model also satisfies all the conclusions of
Chapter~\ref{ch:foundations-v4}: it is orientable, has a finite CW model,
and has $\pi_2(N)=0$ and $\pi_3(N)\cong\mathbb Z$. These are consequences
of its compactness, simple connectedness and dimension, with the
chosen topology throughout. The endpoint below uses simple connectedness;
the other conclusions explain how the smooth model meets the topology
hypotheses used earlier in the geometric argument.

Theorem~\ref{thm:sphere-smooth-endpoint} applies to this smooth model
and produces a diffeomorphism $d:N\cong_{\mathrm{sm}}S^3$.
The topological endpoint map is the prescribed composition
\[
       M\xrightarrow{h}N\xrightarrow{d}S^3,
\]
where $d$ is regarded as a homeomorphism. Our convention that a simply
connected space is nonempty and path connected ensures both hypotheses
needed to construct the smooth model. Thus the construction applies to
every topological manifold in the following statement.

\begin{theorem}[Topological endpoint]
\label{thm:topological-endpoint}
Every compact, Hausdorff, second-countable, simply connected topological
three-manifold is homeomorphic to $S^3$.
\lean{PoincareMT.m77TransportTopologicalHypotheses}
\lean{PoincareMT.m78EndpointTransport}
\lean{PoincareMT.m79TopologicalPoincare}
\lean{PoincareMT.m80EndpointAssembly}
\source{Hamilton1976,Cairns1940,MT2007}
\dcref{Hamilton 1976, Theorem 2(1), p.~64; Cairns 1940, Theorem III, p.~797; Morgan--Tian Corollary 0.2(a)}
\end{theorem}

\begin{proof}
The Hamilton--Cairns construction gives a smooth model $N$ and a specified
homeomorphism $h:M\cong N$. Transport simple connectedness along
$h^{-1}$ and apply Theorem~\ref{thm:sphere-smooth-endpoint} to obtain
$d:N\cong_{\mathrm{sm}}S^3$. Then $d\circ h$ is the required
homeomorphism.
\end{proof}

\section{Dependency of the endpoint}
\label{sec:smoothing-contracts}

The smoothing construction depends on the protected Dehn surfaces,
compact cores, prime replacement and relative rigidity of
Chapters~\ref{ch:dehn-surfaces}--\ref{ch:relative-rigidity}. It is
independent of simple connectedness and of Ricci flow. The smooth sphere
theorem was proved in Chapter~\ref{ch:sphere-reduction}, using the finite
surgery history and extinction. Their combination through the specified
map $d\circ h$ yields the topological conclusion.
